\section{Experiments}
\label{sec:experiments}

\subsection{Genre Classification}

\paragraph{Cross-validation.}
We evaluate both classifiers using stratified 5-fold cross-validation on the full corpus before committing to a final model. Each fold preserves the original class distribution. The mean accuracy across folds is 0.9748 for genre classification and 0.7962 for sentiment classification.

\paragraph{Test set results.}
The final model is trained on the full training partition (80\%, $\approx$1,777 articles) and evaluated on the held-out test set (10\%, 227 articles). It achieves \textbf{98.24\% accuracy}. Per-class precision, recall, and F1 are shown in Table~\ref{tab:classification_report}.

\begin{table}[h]
\centering
\begin{tabular}{lcccc}
\toprule
\textbf{Category} & \textbf{Prec.} & \textbf{Recall} & \textbf{F1} & \textbf{Support} \\
\midrule
Business      & 1.0000 & 0.9608 & 0.9800 & 51 \\
Entertainment & 1.0000 & 0.9750 & 0.9873 & 40 \\
Politics      & 0.9556 & 1.0000 & 0.9773 & 43 \\
Sport         & 0.9811 & 1.0000 & 0.9905 & 52 \\
Tech          & 0.9756 & 0.9756 & 0.9756 & 41 \\
\midrule
\textbf{Weighted avg} & \textbf{0.9829} & \textbf{0.9824} & \textbf{0.9824} & \textbf{227} \\
\bottomrule
\end{tabular}
\caption{Per-class classification report on the 227-article test set.}
\label{tab:classification_report}
\end{table}

The per-class report already shows the remaining errors are small and localized rather than broad failures.

\subsection{Keyword and Entity Extraction}

We present two representative business articles from the corpus to illustrate the keyword and NER pipeline. ``POS-filtered'' means the keyword extractor keeps only nouns and adjectives, which makes the output more readable.

The raw article text uses both the concatenated company name \textit{TimeWarner} and the spaced form \textit{Time Warner}; our tokenizer preserves the concatenated form as the token \textit{timewarner}, while NER still recovers the multiword organization name.

\begin{table*}[t]
\centering
\small
\setlength{\tabcolsep}{4pt}
\begin{tabularx}{\textwidth}{@{}p{2.8cm}>{\raggedright\arraybackslash}X>{\raggedright\arraybackslash}X>{\raggedright\arraybackslash}X@{}}
\toprule
\textbf{Article} & \textbf{TF-IDF} & \textbf{POS-filtered} & \textbf{NER} \\
\midrule
\textit{Time Warner profits} & timewarner, yearearlier, aols, aol, restate, fullyear, 639m, 111bn & timewarner, aol, profits, sales, warner, internet, profit, time & GPE: Ad, Google; PERSON: Time Warner, Warner; ORG: TimeWarner \\
\bottomrule
\addlinespace
\textit{Dollar/Greenspan} & greenspan, 12871, 12974, greenspans, sinche, loosening, halfpoint, deficit & dollar, deficit, months, greenspan, federal, reserve, current & PERSON: Alan Greenspan; ORG: Federal Reserve; FACILITY: White House \\
\bottomrule
\end{tabularx}
\caption{Sample keyword and entity extraction output for two business articles from the corpus.}
\label{tab:samples}
\end{table*}

Several observations stand out. TF-IDF keywords tend to capture proper nouns and rare compound forms that are unique to the article, while part-of-speech filtered keywords tend to be more general and readable. The NER output correctly identifies organizations and people but occasionally misclassifies common nouns as geopolitical entities (e.g., ``Ad'' as a GPE), a known limitation of NLTK's rule-based chunker.

\FloatBarrier

\subsection{Sentiment Analysis}

Table~\ref{tab:sentiment} shows the sentiment predictions for the same three articles from both the NB model and VADER.

\begin{table}[h]
\centering
\begin{tabular}{lccc}
\toprule
\textbf{Article} & \textbf{NB} & \textbf{VADER} & \textbf{Compound} \\
\midrule
Time Warner profits & pos & pos & $+$0.9954 \\
Dollar/Greenspan    & neg & pos & $+$0.9421 \\
Yukos/Rosneft       & neg & neg & $-$0.8860 \\
\bottomrule
\end{tabular}
\caption{NB and VADER sentiment predictions with VADER compound scores for three sample articles. The compound score is VADER's normalized polarity score in $[-1, +1]$; the second article shows a NB/VADER disagreement.}
\label{tab:sentiment}
\end{table}

Agreement between the two systems is mixed. We examine the causes of disagreement in Section~\ref{sec:discussion}.

\FloatBarrier
